\subsection{Interpretation of functions as tempered distributions}

DEFINITION 5. --- A measure \(\nu\) on \(\mathbf{R}^{n}\) is said to be tempered if there exists an integer \(r \in \mathbf{N}\) such that the continuous map \(x \mapsto (1 + \|x\|)^{-r}\) is \(\nu\)-integrable on \(\mathbf{R}^{n}\) .

In other words, a measure \(\nu\) is tempered if there exists \(r \in \mathbf{N}\) such that the function defined by \(x \mapsto \| x\|^{-r}\) is \(\nu\)-integrable on the complement of the unit ball in \(\mathbf{R}^n\) . In particular, every bounded measure on \(\mathbf{R}^n\) is tempered. More generally, if \(f\) is a \(\mu\)-measurable function of polynomial growth and if \(\nu\) is tempered, then the measure \(f \cdot \nu\) is tempered.

The set \(\mathcal{M}^{\mathrm{t}}(\mathbf{R}^{n})\) of tempered measures on \(\mathbf{R}^{n}\) is a vector subspace of the space \(\mathcal{M}(\mathbf{R}^{n};\mathbf{C})\) of complex measures on \(\mathbf{R}^{n}\) .

PROPOSITION 15. --- Let \(\nu\) be a tempered measure on \(\mathbf{R}^{n}\) . The restriction of \(\nu\) to \(\mathcal{S}(\mathbf{R}^{n})\) is a tempered distribution. It is zero if and only if the measure \(\nu\) is zero.

Since \(\nu\) is tempered, there exists a positive integer \(k\) such that the map \(x \mapsto \| x\|^{-k}\) is \(\nu\)-integrable on the complement of the unit ball in \(\mathbf{R}^n\) . For every Schwartz function \(\varphi \in \mathcal{S}(\mathbf{R}^n)\) , we have

\[
| \langle \nu , \varphi \rangle | \leqslant \Bigl (\int_ {\| x \| \leqslant 1} d \nu \Bigr) \widetilde {q} _ {0, 0} (\varphi) + \Bigl (\int_ {\| x \| > 1} \| x \| ^ {- k} d \nu \Bigr) \widetilde {q} _ {0, k} (\varphi),
\]

therefore the restriction of \(\nu\) to \(\mathcal{S}(\mathbf{R}^n)\) is a tempered distribution.

The last assertion follows from prop. 6 of IV, p. 206, since \(\mathcal{D}(\mathbf{R}^n)\) is contained in \(\mathcal{S}(\mathbf{R}^n)\) .

We will identify the space \(\mathcal{M}^{\mathrm{t}}(\mathbf{R}^{n})\) with a subspace of \(\mathcal{S}'(\mathbf{R}^{n})\) .

PROPOSITION 16. --- Let \(p \in [1, +\infty]\) and \(f \in \mathcal{L}^{p}(\mathbf{R}^{n})\) . Then the measure \(f \cdot \mu\) with density \(f\) with respect to Lebesgue measure is tempered. The map \(f \mapsto f \cdot \mu\) from \(\mathrm{L}^{p}(\mathbf{R}^{n})\) into \(\mathcal{S}'(\mathbf{R}^{n})\) thus defined is continuous.

Let \(q\) be the conjugate exponent of \(p\) and let \(r \geqslant 0\) be such that \(rq > n\) . For every \(x \in \mathbf{R}^n\) , write \(g(x) = (1 + \|x\|)^{-r}\) . The function \(g\) belongs to \(\mathcal{L}^q(\mathbf{R}^n)\) by prop. 3 of IV, p. 199. By Hölder's inequality, we have

\[
\int_ {\mathbf {R} ^ {n}} ^ {*} (1 + \| x \|) ^ {- r} | f (x) | d \mu (x) \leqslant \mathrm{N} _ {q} (g) \mathrm{N} _ {p} (f) <   + \infty ,
\]

therefore the measure \(f \cdot \mu\) is tempered.

Let \(f \in \mathrm{L}^p(\mathbf{R}^n)\) and \(\varphi \in \mathcal{S}(\mathbf{R}^n)\) . Hölder's inequality implies

\[
| \langle f \cdot \mu , \varphi \rangle | = \left| \int_ {\mathbf {R} ^ {n}} f (x) \varphi (x) d \mu (x) \right| \leqslant \| f \| _ {p} \| \varphi \| _ {q}
\]

and the continuity of the mapping \(f \mapsto f \cdot \mu\) then follows from prop. 13 of IV, p. 213.

For every \(p \in [1, +\infty]\) , we shall identify \(\mathrm{L}^p(\mathbf{R}^n)\) with a subspace of \(\mathcal{S}'(\mathbf{R}^n)\) by the linear mapping \(f \mapsto f \cdot \mu\) .

PROPOSITION 17. --- The spaces \(\mathcal{D}(\mathbf{R}^{n})\) and \(\mathcal{S}(\mathbf{R}^{n})\) are dense in \(\mathcal{S}'(\mathbf{R}^{n})\) .

It suffices to show that \(\mathcal{D}(\mathbf{R}^{n})\) is dense in \(\mathcal{S}'(\mathbf{R}^{n})\) . Let \(\lambda\) be a continuous linear form on \(\mathcal{S}'(\mathbf{R}^{n})\) vanishing on \(\mathcal{D}(\mathbf{R}^{n})\) . Since the space \(\mathcal{S}(\mathbf{R}^{n})\) is reflexive, there exists a function \(\varphi \in \mathcal{S}(\mathbf{R}^{n})\) such that \(\lambda(f) = \langle f, \varphi \rangle\) for all \(f \in \mathcal{S}'(\mathbf{R}^{n})\) . We therefore have

\[
0 = \lambda (\psi) = \langle \psi , \varphi \rangle = \int_ {\mathbf {R} ^ {n}} \psi \varphi d \mu
\]

for all \(\psi\in\mathcal{D}(\mathbf{R}^{n})\) . The measure \(\varphi\cdot\mu\) on \(\mathbf{R}^{n}\) is therefore zero (prop. 6 of IV, p. 206), whence \(\varphi=0\) . The proposition then follows from the Hahn–Banach theorem (EVT, II, p. 46, cor. 1).

